% ============ 第4章 参考文献(原书 pp.378-380) ============
\section*{参考文献}\addcontentsline{toc}{section}{参考文献}
\begin{thebibliography}{99}
% 条目 [13],[14] 为日文原条目,其中个别汉字(如 掛,ー 等)为 Fandol 字体所无,
% 此处临时切换至系统日文字体渲染(条目前后不改变).
\newCJKfontfamily\japonfont{MS Gothic}
\begin{sloppypar}
\item[] 关于用驻相法处理振荡积分算子以及 Fourier 限制问题(Tomas--Stein 定理)的部分,依照 \cite{ch4:1} 的最初部分整理而成.作为这一领域的基本教材,也列出 \cite{ch4:2}.
\bibitem[1]{ch4:1} C. D. Sogge, Fourier Integrals in Classical Analysis, Cambridge Tracts in Math. 105 (1993).
\bibitem[2]{ch4:2} E. M. Stein, Harmonic Analysis: Real-Variable Methods, Orthogonality, and Oscillatory Integrals, Princeton University Press, 1993.
\item[] 关于 Nikodym 极大函数(Wolff 定理)的部分,基本上依据 Wolff 的原始论文 \cite{ch4:3},但按照将其整理改编后的 \cite{ch4:4} 编写.
\bibitem[3]{ch4:3} T. Wolff, An improved bound for Kakeya type maximal functions, Rev. Mat. Iberoamericana 11 (1995), 651--674.
\bibitem[4]{ch4:4} H. Tanaka, The Fefferman-Stein type inequality for the Kakeya maximal operator in Wolff's range, Proc. Amer. Math. Soc. 133 (2005), 763--772.
\item[] 作为与 Bochner--Riesz 平均和 Nikodym 极大函数部分相关的论文,列出如下.
\bibitem[5]{ch4:5} A. C\'ordoba, The Kakeya maximal function and the spherical summation multiplier, Amer. J. Math. 99 (1977), 1-22.
\bibitem[6]{ch4:6} A. C\'ordoba, A note on Bochner-Riesz operators, Duke Math. J. 46 (1979), 505--511.
\bibitem[7]{ch4:7} J. Bourgain, Besicovitch type maximal operators and applications to Fourier analysis, Geom. Funct. Anal. 1 (1991), no. 2, 147--187.
\item[] 作为关于 Fourier 限制问题的综述,列出如下.
\bibitem[8]{ch4:8} T. Tao, Recent progress on the Restriction conjecture, Fourier analysis and convexity, 217--243, Appl. Numer. Harmon. Anal., 2004.
\item[] 作为关于 Kakeya 问题的综述,列出如下.
\bibitem[9]{ch4:9} J. Bourgain, Harmonic analysis and combinatorics: How much may they contribute to each other?, Mathematics: Frontiers and perspectives, IMU/Amer. Math. Society, 2000, 13--32.
\bibitem[10]{ch4:10} T. Wolff, Recent work connected with the Kakeya problem, Prospects in mathematics (Princeton, NJ, 1996), 129--162, Amer. Math. Soc., Providence, RI, 1999.
\bibitem[11]{ch4:11} T. Tao, From rotating needles to stability of waves: emerging connections between combinatorics, analysis, and PDE, Notices Amer. Math. Soc. 48 (2001), no. 3, 294--303.
\bibitem[12]{ch4:12} N. H. Katz and T. Tao, Recent progress on the Kakeya conjecture, Publ. Mat. 2002, Vol. Extra, 161--179.
\bibitem[13]{ch4:13} {\japonfont 田中仁, 掛谷問題について,} 数学 57 (2005), no. 2, 113--129.
\item[] 作为论述 Kakeya 集构造的日语教材,列出如下.
\bibitem[14]{ch4:14} {\japonfont 新井仁之, ルベーグ積分講義, 日本評論社,} 2003.
\item[] 作为用组合论方法研究 Nikodym 极大函数所得的结果,列出如下论文.
\bibitem[15]{ch4:15} J. Bourgain, On the dimension of Kakeya sets and related maximal inequalities, Geom. Funct. Anal. 9 (1999), no. 2, 256--282.
\bibitem[16]{ch4:16} N. H. Katz and T. Tao, Bounds on arithmetic projections, and applications to the Kakeya conjecture, Math. Res. Lett. 6 (1999), 625--630.
\bibitem[17]{ch4:17} N. H. Katz and T. Tao, New bounds for Kakeya problems, J. Anal. Math. 87 (2002), 221--263.
\item[] 作为用双线性方法研究 Fourier 限制问题以及 Bochner--Riesz 平均问题所得的结果,列出如下论文.
\bibitem[18]{ch4:18} T. Tao, A. Vargas and L. Vega, A bilinear approach to the restriction and Kakeya conjectures, J. Amer. Math. Soc. 11 (1998), 967--1000.
\bibitem[19]{ch4:19} T. Wolff, A sharp bilinear cone restriction estimate, Annals of Math. 153 (2001), 661-698.
\bibitem[20]{ch4:20} T. Tao, Endpoint bilinear restriction theorems for the cone, and some sharp null form estimates, Math. Z. 238 (2001), no. 2, 215-268.
\bibitem[21]{ch4:21} T. Tao, A sharp bilinear restriction estimate on paraboloids, Geom. Funct. Anal. 13 (2003), 1359--1384.
\bibitem[22]{ch4:22} S. Lee, Improved bounds for Bochner-Riesz and maximal Bochner-Riesz operators, Duke Math. J. 122 (2004), no. 1, 205-232.
\bibitem[23]{ch4:23} S. Lee, Bilinear restriction estimates for surfaces with curvatures of different signs. Trans. Amer. Math. Soc. 358 (2006), no. 8, 3511-3533.
\end{sloppypar}
\end{thebibliography}